% !TeX spellcheck = it_IT

\documentclass[]{beamer}

\usepackage[utf8]{inputenc}
\usepackage[italian]{babel}
\usepackage{multicol}
\usetheme{metropolis}

\usepackage{listings}

\usepackage{tikz}
\usetikzlibrary{shapes,arrows,positioning}

\tikzstyle{startstop} = [ellipse, draw, fill=bg, text=structure, text width=3em, text badly centered, inner sep=0pt]
\tikzstyle{block} = [rectangle, draw, fill=bg, text=structure, text width=5em, text centered]
\tikzstyle{decision} = [diamond, aspect=2, draw, fill=bg, text=structure, text width=3em, text badly centered, inner sep=0pt]
\tikzstyle{input} = [rectangle, draw, fill=bg, text=structure, text width=4em, text centered, rounded rectangle, rounded rectangle east arc=none]
\tikzstyle{output} = [rectangle, draw, fill=bg, text=structure, text width=4em, text centered, rounded rectangle, rounded rectangle west arc=none]

\tikzstyle{line} = [draw, -latex']


\title{Informatica B - Esercitazione 5}
\subtitle{Introduzione a Matlab}

\author{Stefano Cereda\\
		stefano1.cereda@mail.polimi.it
	}
\date{31/10/2017}
\institute[PoliMi]{\vspace{0.5cm}\centering Politecnico di Milano \\ \vspace{0.2cm}
	\includegraphics[width=0.2\linewidth]{./logopolimi}}
%\logo{\includegraphics[width=15mm]{./logopolimi}}

\setbeamercovered{invisible}

\makeindex

\begin{document}
\begin{frame}
	\maketitle
\end{frame}




\begin{frame}{Esercizio 2 - costrutto if in Matlab: equazioni di secondo grado}
Scaricare e installare matlab

Script: commenti(nome e descrizione; data e autori; def variabili), input, calcolo, output.

Si chieda l’inserimento dei tre coefficienti di una equazione di secondo grado: $ax^2+bx+c=0$ e si stampi il valore delle radici reali.
\end{frame}

\begin{frame}{Esercizio 3 - grafici in Matlab}
Si chieda l’inserimento dei tre coefficienti di una equazione di secondo grado: $ax^2+bx+c=0$ e si dia una rappresentazione della curva.
\end{frame}

\begin{frame}{Esercizio 4 - costrutto while in Matlab: fattoriale}
Si scriva uno script che chieda l'inserimento di un numero $n$ e ne calcoli e visualizzi il fattoriale con un ciclo while.
\end{frame}

\begin{frame}{Esercizio 5 - Matlab: matrice triangolare}
Sc scriva una script in Matlab/Octave che svolge le seguenti operazioni:
\begin{enumerate}
	\item Definisce una matrice quadrata di interi di dimensione $N\times N$, con $N=5$.
	\item Acquisisce da tastiera una sequenza di valori interi che memorizza nella	porzione triangolare superiore della matrice. La parte triangolare superiore della matrice è costituita dalle celle poste sulla diagonale principale e al di sopra di tale diagonale.
	\item Copia in un array di dimensioni opportune tutti i valori della porzione triangolare di matrice che sono maggiori di -10.
	\item Stampa il contenuto dell’array.
\end{enumerate}
\end{frame}

\end{document}